\section{Industrielle KI-Bildanalyse als methodisches Vorbild}
\label{sec:grundlagen}

\subsection{Technische Grundlagen und industrielle Anwendung} \label{sec:cnns}

Convolutional Neural Networks (CNNs) bilden das Fundament moderner KI-Bildanalyse in
Industrie und Medizin gleichermaßen. Durch hierarchische Faltungsoperationen extrahieren
sie räumliche Merkmale: Frühe Schichten erfassen Kanten und Texturen, tiefere Schichten
kombinieren diese zu semantischen Strukturen -- etwa Defektmuster oder pathologische
Veränderungen. Transfer Learning ergänzt diesen Ansatz: Ein auf einem großen allgemeinen
Datensatz vortrainiertes Modell wird durch Fine-Tuning auf einem domänenspezifischen
Zieldatensatz spezialisiert. Da sowohl industrielle als auch medizinische Bilddatensätze
annotationsaufwändig und vergleichsweise klein sind, stellt dieser Ansatz in beiden
Domänen die entscheidende Voraussetzung für Praxistauglichkeit dar \cite{transferMedical2024}.

In der industriellen Qualitätskontrolle sind CNN-basierte Systeme produktiv etabliert.
Systeme zur automatisierten Erkennung von Materialfehlern auf Baustoffen kombinieren
Bildsegmentierung zur Lokalisation des Defektbereichs mit Klassifikation zur Einordnung
des Fehlertyps (z.\,B.\ Risse, Ablösungen, Fremdeinschlüsse). Die Ausgabe ist ein
strukturierter Regelentscheid in Echtzeit, ohne menschliche Zwischeninstanz. Die
zentralen Systemeigenschaften -- Skalierbarkeit, Reproduzierbarkeit und Unabhängigkeit
von individueller Ermüdung -- sind genau jene, die der radiologischen Befundung unter
wachsendem Ressourcendruck fehlen (vgl. Abschnitt~\ref{sec:problemstellung}).

\subsection{Struktureller Vergleich beider Domänen} \label{sec:parallelen}

Tabelle~\ref{tab:vergleich} zeigt, dass industrielle Qualitätskontrolle und medizinische
Bilddiagnostik strukturell identische Aufgaben lösen.

\begin{table}[htbp]
    \centering
    \caption{Struktureller Vergleich: Industrielle CV vs.\ medizinische Bildanalyse}
    \label{tab:vergleich}
    \begin{tabularx}{\textwidth}{|l|X|X|}
        \hline
        \textbf{Aspekt} & \textbf{Industrielle CV} & \textbf{Medizinische Bildanalyse} \\
        \hline
        Eingabedaten    & Kameraaufnahmen (Baustoffe, Produkte) & Röntgen, CT, MRT \\
        Anomalieklassen & Risse, Ablösungen, Einschlüsse        & Tumore, Frakturen, Läsionen \\
        Kernmethode     & Segmentierung + Klassifikation         & Segmentierung + Klassifikation \\
        Ausgabe         & Defekttyp + Regelentscheid             & Befund + Diagnoseempfehlung \\
        Betriebsstatus  & produktiv, skaliert, erprobt           & in Einführung und Erprobung \\
        \hline
    \end{tabularx}
    \source{Eigene Darstellung}
\end{table}

Die nachfolgenden Abschnitte untersuchen, wie dieser Transfer empirisch validiert wird
und welche infrastrukturellen Voraussetzungen er im Gesundheitswesen erfordert.
